\chapter{Koneksje i transport równoległy}
\label{ch:koneksje}
\index{koneksja}

Wektor styczny w punkcie $p$ należy do $T_pM$, a wektor
w punkcie $q$ do innej przestrzeni $T_qM$. Samo odejmowanie
tych wektorów nie ma więc znaczenia. Koneksja mówi,
jak różniczkować pole wektorowe, gdy przemieszczamy się
po rozmaitości. W metryce riemannowskiej szczególna
koneksja pozwala przenosić wektory bez zmiany ich
iloczynów skalarnych. Będziemy używać definicji
pochodnej pola z Rozdziału~\ref{ch:algebra-tensorowa}
oraz metryk z Rozdziału~\ref{ch:metryka-riemannowska}.

\section{Dlaczego zwykła pochodna współrzędnych nie wystarcza?}
\label{sec:koneksja-motywacja}

Na płaszczyźnie euklidesowej wektor można przesuwać
równolegle, bo wszystkie przestrzenie styczne mają
kanoniczny opis przez $\R^2$. Nawet tu pochodne
\emph{współczynników} zależą od bazy: wektor stale
wskazujący na wschód ma stałe współczynniki w bazie
kartezjańskiej, lecz zmienne w obracającej się bazie
biegunowej. Na sferze potrzebny jest ponadto przepis,
jak porównać wektory styczne w jej różnych punktach.

\begin{figure}[htbp]
\centering
\begin{tikzpicture}[scale=0.95,>=Latex]
  \fill[manifoldblue!8] (0,0) ellipse (3.25 and 1.27);
  \draw[manifoldblue!65,thick] (0,0) ellipse (3.25 and 1.27);
  \draw[manifoldblue,very thick,->] (-2.5,-.32) .. controls (-1.3,.83) and (.36,.64) .. (2.30,.31);
  \path (-2.5,-.32) .. controls (-1.3,.83) and (.36,.64) .. (2.30,.31)
       coordinate[pos=.22] (pa) coordinate[pos=.85] (qa);
  \fill[manifoldblue] (pa) circle (2.2pt) node[below=6pt] {$p$};
  \fill[manifoldblue] (qa) circle (2.2pt) node[below=6pt] {$q$};
  \draw[manifoldred,ultra thick,->] (pa) -- ++(.35,.73) node[above] {$v\in T_pM$};
  \draw[manifoldred,ultra thick,->] (qa) -- ++(.54,.62) node[above] {$P_\gamma v\in T_qM$};
  \node[manifoldblue] at (0,-.62) {$\gamma$};
  \node[align=center,text width=7.9cm,below] at (0,-1.36)
    {Koneksja określa transport wzdłuż $\gamma$; wynik może zależeć od drogi.};
\end{tikzpicture}
\caption{Schemat transportu wzdłuż drogi. Strzałki reprezentują wektory w różnych przestrzeniach stycznych; ich położenie na kartce nie wyznacza metryki ani wielkości obrotu.}
\label{fig:koneksja-transport-idea}
\end{figure}

\begin{definicja}[Koneksja afiniczna]
\index{koneksja afiniczna@Koneksja afiniczna}
\label{def:koneksja-afiniczna}
Koneksją afiniczną na $TM$ nazywamy przyporządkowanie
$(X,Y)\mapsto\nabla_XY$ pól wektorowych, liniowe
nad $\R$ w obu argumentach, takie że dla $f\in C^\infty(M)$
\begin{equation}\label{eq:koneksja-aksjomaty}
 \nabla_{fX}Y=f\nabla_XY,
 \qquad \nabla_X(fY)=X(f)Y+f\nabla_XY.
\end{equation}
Pierwszy wzór pozwoli wykazać, że $\nabla_XY$ w $p$ zależy
tylko od wektora $X_p$. Drugi jest regułą iloczynu:
zmienia się zarówno wielkość współczynnika $f$,
jak i samo pole $Y$.
\end{definicja}

\begin{lemat}[Lokalność koneksji]
\label{lem:koneksja-lokalnosc}
Wartość $(\nabla_XY)_p$ zależy tylko od ograniczeń $X,Y$
do dowolnego otoczenia $p$. Można więc stosować koneksję
do pól określonych jedynie na otwartym podzbiorze $M$.
W pierwszym argumencie wystarcza sama wartość $X_p$.
\end{lemat}
\begin{proof}
Jeśli pole $Z$ znika na otoczeniu $U$ punktu $p$, wybieramy
gładką funkcję $\chi$ o nośniku zawartym w $U$, równą $1$
w pobliżu $p$. Takie funkcje daje konstrukcja z
Sekcji~\ref{sec:formy-podzial-jednosci}. Ponieważ $\chi Z=0$,
aksjomaty koneksji dają
\[
 0=\nabla_{\chi Z}Y=\chi\nabla_ZY,\qquad
 0=\nabla_X(\chi Z)=X(\chi)Z+\chi\nabla_XZ.
\]
Po obliczeniu w $p$ otrzymujemy $(\nabla_ZY)_p=0$
i $(\nabla_XZ)_p=0$. Odejmowanie pól dowodzi lokalności.
Pole określone na $U$ mnożymy przez taką funkcję $\chi$
i przedłużamy zerem poza $U$; wynik obliczenia w $p$
nie zależy od przedłużenia. Te lokalne definicje sklejają się.
W lokalnej ramce $X=X^i\partial_i$ pierwsza reguła daje
$(\nabla_XY)_p=X^i(p)(\nabla_{\partial_i}Y)_p$,
co dowodzi ostatniego zdania. Drugi argument wymaga
także pierwszych pochodnych współczynników pola $Y$.
\end{proof}

\begin{stwierdzenie}[Współczynniki lokalne]
\label{prop:koneksja-lokalna}
W mapie $(x^1,\ldots,x^n)$ koneksję wyznaczają
funkcje $\Gamma^k_{ij}$, zwane symbolami Christoffela,
określone przez
\begin{equation}\label{eq:def-christoffel}
 \nabla_{\partial_i}\partial_j
   =\sum_k\Gamma^k_{ij}\partial_k.
\end{equation}
Dla $X=X^i\partial_i$ i $Y=Y^j\partial_j$ mamy
\begin{equation}\label{eq:koneksja-wspolrzedne}
 (\nabla_XY)^k=X^i\bigl(\partial_iY^k+
                         \Gamma^k_{ij}Y^j\bigr).
\end{equation}
W powtarzających się indeksach stosujemy sumowanie.
\end{stwierdzenie}
\begin{proof}
Z pierwszego aksjomatu
$\nabla_XY=X^i\nabla_{\partial_i}(Y^j\partial_j)$.
Z drugiego
$\nabla_{\partial_i}(Y^j\partial_j)
=(\partial_iY^j)\partial_j
 +Y^j\Gamma^k_{ij}\partial_k$.
Zbierając współczynniki przy $\partial_k$,
otrzymujemy podany wzór. I odwrotnie:
gładkie funkcje $\Gamma^k_{ij}$ w mapie definiują
tam koneksję przez ten wzór.
\end{proof}

\begin{uwaga}[Symbole nie są tensorem]
\label{uw:symbole-nie-tensor}
Przy zamianie współrzędnych $x=x(y)$ stosujemy
$\partial_{y^a}=(\partial x^i/\partial y^a)\partial_{x^i}$
oraz regułę iloczynu. Jeśli $B^i_a=\partial x^i/\partial y^a$,
to przed zmianą bazy w wyniku mamy
\[
 \nabla_{\partial_{y^a}}\partial_{y^b}
 =\bigl(B^i_a\partial_{x^i}B^k_b
       +B^i_aB^j_b\Gamma^k_{ij}\bigr)\partial_{x^k}.
\]
Pierwszy składnik to $\partial^2x^k/\partial y^a\partial y^b$.
Po zapisaniu $\partial_{x^k}=(\partial y^c/\partial x^k)
\partial_{y^c}$ otrzymujemy
\begin{equation}\label{eq:transformacja-christoffel}
 \widetilde\Gamma^c_{ab}
 =\frac{\partial y^c}{\partial x^k}
  \left(
   \frac{\partial x^i}{\partial y^a}
   \frac{\partial x^j}{\partial y^b}\Gamma^k_{ij}
   +\frac{\partial^2x^k}{\partial y^a\partial y^b}
  \right).
\end{equation}
Druga pochodna jest dodatkowym składnikiem;
same $\Gamma^k_{ij}$ nie tworzą tensora.
To pozwala mieć zerowe symbole w kartezjańskim
opisie płaszczyzny i niezerowe w biegunowym,
choć geometrycznie jest to ta sama płaszczyzna.
\end{uwaga}

\begin{definicja}[Torsja i zgodność z metryką]
\index{torsja i zgodnosc z metryka@Torsja i zgodność z metryką}
\label{def:torsja-koneksja}
Torsją koneksji jest pole tensorowe
\[
 T(X,Y)=\nabla_XY-\nabla_YX-[X,Y].
\]
Koneksja jest \emph{bez torsji}, gdy $T=0$.
Jeżeli $g$ jest metryką, koneksja jest z nią
\emph{zgodna}, gdy dla dowolnych pól
\begin{equation}\label{eq:zgodnosc-metryczna}
 X\bigl(g(Y,Z)\bigr)
 =g(\nabla_XY,Z)+g(Y,\nabla_XZ).
\end{equation}
\end{definicja}

\begin{stwierdzenie}[Co zależy tylko od wartości w punkcie?]
\label{prop:torsja-tensor-roznica}
Torsja jest tensorem typu $(1,2)$. Jeżeli
$\nabla$ i $\widehat\nabla$ są koneksjami,
to $A(X,Y)=\widehat\nabla_XY-\nabla_XY$
również jest tensorem typu $(1,2)$.
\end{stwierdzenie}
\begin{proof}
W pierwszym argumencie dla torsji używamy
$[fX,Y]=f[X,Y]-Y(f)X$. Po rozwinięciu
\[
 T(fX,Y)=f\nabla_XY-Y(f)X-f\nabla_YX
          -f[X,Y]+Y(f)X=fT(X,Y).
\]
Analogicznie $[X,fY]=f[X,Y]+X(f)Y$,
więc wyrazy $X(f)Y$ w $T(X,fY)$ się skracają.
Ponieważ $\nabla$ i $\widehat\nabla$ mają tę
samą regułę iloczynu, również w ich różnicy
kasują się składniki $X(f)Y$.
To wyjaśnia, czemu przestrzeń wszystkich koneksji
jest afiniczna: do jednej koneksji można dodać
dowolny tensor takiego typu. W mapie
$T^k_{ij}=\Gamma^k_{ij}-\Gamma^k_{ji}$,
bo $[\partial_i,\partial_j]=0$.
\end{proof}

\section{Koneksja Leviego-Civity i wzór Koszula}
\label{sec:levi-civita-koszul}

Metryka dostarcza iloczynu skalarnego w każdym
punkcie. Wybieramy koneksję, która zachowuje ten
iloczyn podczas przemieszczania się, i żądamy
zerowej torsji. Te dwa warunki usuwają dowolność.

\begin{twierdzenie}[Istnienie i jednoznaczność koneksji Leviego-Civity]
\label{thm:levi-civita}
Każda metryka riemannowska $g$ na $M$ wyznacza
dokładnie jedną koneksję $\nabla$ zgodną z $g$
i bez torsji. Ta sama teza zachodzi dla metryki
pseudoriemannowskiej, w szczególności
Schwarzschilda i Kerra.
\end{twierdzenie}
\begin{proof}
\emph{Krok 1: wyprowadzenie wzoru i jednoznaczność.}
Załóżmy, że taka koneksja istnieje. Zapiszmy
\eqref{eq:zgodnosc-metryczna} kolejno dla
$(X,Y,Z)$, $(Y,Z,X)$ oraz $(Z,X,Y)$:
\begin{align*}
 Xg(Y,Z)&=g(\nabla_XY,Z)+g(Y,\nabla_XZ),\\
 Yg(Z,X)&=g(\nabla_YZ,X)+g(Z,\nabla_YX),\\
 Zg(X,Y)&=g(\nabla_ZX,Y)+g(X,\nabla_ZY).
\end{align*}
Dodajemy dwa pierwsze równania i odejmujemy
trzecie. Brak torsji daje
$\nabla_YX=\nabla_XY-[X,Y]$,
$\nabla_ZX=\nabla_XZ-[X,Z]$ oraz
$\nabla_ZY=\nabla_YZ-[Y,Z]$.
Po podstawieniu pary z $\nabla_XZ$ i pary
z $\nabla_YZ$ się skracają. Zostaje
\begin{align}\label{eq:koszul}
2g(\nabla_XY,Z)={}&Xg(Y,Z)+Yg(Z,X)-Zg(X,Y)\notag\\
 &+g([X,Y],Z)-g([X,Z],Y)-g([Y,Z],X).
\end{align}
Jest to \emph{wzór Koszula}. Jego prawa strona
zależy wyłącznie od $g,X,Y,Z$. Niezdegenerowanie
$g_p$ sprawia, że wartości $g_p(W,Z_p)$ dla
wszystkich $Z_p$ wyznaczają jedyny wektor $W$.
Stąd co najwyżej jedna taka koneksja.

\emph{Krok 2: konstrukcja.} Oznaczmy prawą stronę
\eqref{eq:koszul} przez $K(X,Y,Z)$ i zdefiniujmy
$\nabla_XY$ przez $2g(\nabla_XY,Z)=K(X,Y,Z)$.
Najpierw trzeba sprawdzić, że prawa strona
jest $C^\infty$-liniowa względem $Z$.
Po zastąpieniu $Z$ przez $fZ$ składniki z
$X(f)g(Y,Z)$ pochodzące od
$Xg(Y,fZ)$ i $-g([X,fZ],Y)$ mają przeciwne znaki;
tak samo znoszą się $Y(f)g(X,Z)$ z
$Yg(fZ,X)$ i $-g([Y,fZ],X)$.
Reszta jest pomnożona przez $f$.
Zatem $K(X,Y,\cdot)$ jest gładką $1$-formą,
której współczynniki w mapie to $K(X,Y,\partial_i)$.
Jawnie
\[
 (\nabla_XY)^j=\tfrac12g^{ji}K(X,Y,\partial_i).
\]
Macierz $G^{-1}$ jest gładka, więc otrzymujemy gładkie pole.
Definicje w różnych mapach są zgodne, ponieważ wszystkie
spełniają tę samą równość z $g$, a niezdegenerowanie
metryki zapewnia punktową jednoznaczność.

Sprawdźmy aksjomaty koneksji. W $K(fX,Y,Z)$
dodatkowe $Y(f)g(X,Z)$ z $Yg(Z,fX)$
znosi $-Y(f)g(X,Z)$ z $g([fX,Y],Z)$.
Dodatkowe $-Z(f)g(X,Y)$ z $-Zg(fX,Y)$
znosi $+Z(f)g(X,Y)$ z $-g([fX,Z],Y)$.
Wobec tego $K(fX,Y,Z)=fK(X,Y,Z)$.
W $K(X,fY,Z)$ dwukrotnie pojawia się
$X(f)g(Y,Z)$: z $Xg(fY,Z)$ oraz
$g([X,fY],Z)$. Wyrazy $Z(f)g(X,Y)$
z $-Zg(X,fY)$ i $-g([fY,Z],X)$
znoszą się. Dostajemy
\[
 K(X,fY,Z)=fK(X,Y,Z)+2X(f)g(Y,Z).
\]
Po podzieleniu przez $2$ i użyciu
niezdegenerowania $g$ są to dokładnie
aksjomaty \eqref{eq:koneksja-aksjomaty}.

\emph{Krok 3: sprawdzenie warunków.}
Różnica $K(X,Y,Z)-K(Y,X,Z)$ jest równa
$2g([X,Y],Z)$: pochodne metryki
się kasują, a antysymetria nawiasu daje
wyraz podwojony. Stąd
$\nabla_XY-\nabla_YX=[X,Y]$.
Dodając $K(X,Y,Z)$ i $K(X,Z,Y)$,
otrzymujemy $2Xg(Y,Z)$: pozostałe pochodne
i wyrazy z nawiasami redukują się parami.
To jest \eqref{eq:zgodnosc-metryczna}.
Użyliśmy jedynie niezdegenerowania macierzy $G$,
nie jej dodatniej określoności; dowód dotyczy
również metryk lorentzowskich.
\end{proof}

W szczególności istnienie metryki riemannowskiej
z Twierdzenia~\ref{thm:istnienie-metryki}
zapewnia istnienie \emph{jakiejś} koneksji
afinicznej na każdej rozmaitości z naszych założeń.
Koneksji jest zwykle wiele, ale dla ustalonego $g$
warunki z twierdzenia wybierają jedną.

\begin{stwierdzenie}[Podrozmaitość w przestrzeni euklidesowej]
\label{prop:koneksja-projekcja}
Jeżeli $M\subset\R^N$ ma metrykę indukowaną,
to jej koneksja Leviego-Civity jest równa
\[
 \nabla_XY=\operatorname{pr}_{TM}(D_XY),
\]
gdzie $D_XY$ oznacza zwykłą pochodną pola $Y$
traktowanego jako odwzorowanie do $\R^N$.
\end{stwierdzenie}
\begin{proof}
Rzut ortogonalny jest liniowy w punkcie,
a pochodna euklidesowa spełnia regułę iloczynu,
więc wzór określa koneksję. Pochodna iloczynu
skalarnego w $\R^N$ daje
$X\langle Y,Z\rangle
=\langle D_XY,Z\rangle+\langle Y,D_XZ\rangle$.
Składowe normalne obu pochodnych są prostopadłe
do $Y$ lub $Z$ i znikają z tej równości;
koneksja jest zgodna z metryką indukowaną.
Ponadto w lokalnej parametryzacji mieszane
pochodne drugiego rzędu komutują, więc
$D_XY-D_YX=[X,Y]$; po rzucie otrzymujemy
zerową torsję. Stosujemy Twierdzenie~\ref{thm:levi-civita}.
\end{proof}

\section{Symbole Christoffela w geometriach modelowych}
\label{sec:christoffel-przyklady}

W mapie pola bazowe komutują, więc wzór Koszula
z $X=\partial_i$, $Y=\partial_j$, $Z=\partial_\ell$
przyjmuje postać
\[
 2g_{k\ell}\Gamma^k_{ij}
 =\partial_i g_{j\ell}+\partial_j g_{i\ell}
    -\partial_\ell g_{ij}.
\]
Mnożymy przez $g^{m\ell}$, gdzie
$(g^{m\ell})=G^{-1}$. Ponieważ
$g^{m\ell}g_{k\ell}=\delta^m_k$, otrzymujemy
\begin{equation}\label{eq:christoffel-metryka}
 \boxed{\displaystyle
 \Gamma^m_{ij}=\frac12 g^{m\ell}
  (\partial_i g_{j\ell}+\partial_j g_{i\ell}
                         -\partial_\ell g_{ij}).}
\end{equation}
Algorytm jest prosty: zapisać $G$, odwrócić ją,
obliczyć pochodne współczynników, podstawić.
Duże metryki nie wymagają zapamiętania wielkiej
listy symboli; zwykle liczymy tylko te, które
potrzebne są w danym równaniu.

\begin{przyklad}[Płaszczyzna euklidesowa w dwóch mapach]
\label{prz:christoffel-biegunowe}
W kartezjańskiej mapie $g=\dd x^2+\dd y^2$,
wszystkie $g_{ij}$ są stałe, zatem $\Gamma^k_{ij}=0$.
W biegunowej mapie $x=r\cos\phi$, $y=r\sin\phi$
(dla $r>0$ i $\phi$ w otwartym przedziale długości co najwyżej
$2\pi$) ta sama metryka jest
$g=\dd r^2+r^2\dd\phi^2$.
Macierze $G=\operatorname{diag}(1,r^2)$
i $G^{-1}=\operatorname{diag}(1,r^{-2})$ oraz
$\partial_r g_{\phi\phi}=2r$ dają
\begin{equation}\label{eq:gamma-biegunowe}
 \Gamma^r_{\phi\phi}=-r,
 \qquad \Gamma^\phi_{r\phi}
 =\Gamma^\phi_{\phi r}=\frac1r,
\end{equation}
a wszystkie pozostałe symbole są zerowe.
Na przykład
$\nabla_{\partial_\phi}\partial_\phi=-r\partial_r$:
wektor $\partial_\phi$ w miarę obrotu układu
zmienia kierunek nawet wtedy, gdy płaszczyzna
pozostaje płaska. Wzór $\nabla_XY$ uwzględnia
ten obrót bazy.
\end{przyklad}

\begin{przyklad}[Sfera o promieniu $R$]
\label{prz:christoffel-sfera}
Z \eqref{eq:metryka-sfery} odczytujemy
$G=\operatorname{diag}(R^2,R^2\sin^2\theta)$.
Jedyna niezerowa pochodna współczynników to
$\partial_\theta g_{\phi\phi}
=2R^2\sin\theta\cos\theta$. Stąd
\begin{equation}\label{eq:gamma-sfera}
 \Gamma^\theta_{\phi\phi}
   =-\sin\theta\cos\theta,
 \qquad
 \Gamma^\phi_{\theta\phi}
 =\Gamma^\phi_{\phi\theta}=\cot\theta.
\end{equation}
Brak $R$ w tych symbolach jest zgodny z tym,
że stałe przemnożenie metryki nie zmienia
jej koneksji. Bieguny wymagają innej mapy;
$\cot\theta$ nie opisuje tam osobliwości sfery.
\end{przyklad}

\begin{przyklad}[Górna półprzestrzeń hiperboliczna]
\label{prz:christoffel-hiperboliczna}
Dla dwuwymiarowego modelu
$g=y^{-2}(\dd x^2+\dd y^2)$, $y>0$,
macierze $G=y^{-2}I$, $G^{-1}=y^2I$.
Ponieważ $\partial_y(y^{-2})=-2y^{-3}$,
\eqref{eq:christoffel-metryka} daje
\begin{equation}\label{eq:gamma-hiperboliczna}
 \Gamma^y_{xx}=\frac1y,\quad
 \Gamma^y_{yy}=-\frac1y,\quad
 \Gamma^x_{xy}=\Gamma^x_{yx}=-\frac1y.
\end{equation}
W wyższych wymiarach, dla $g=y^{-2}
(\sum_a\dd x_a^2+\dd y^2)$, podobnie
$\Gamma^y_{x_ax_b}=\delta_{ab}/y$,
$\Gamma^{x_a}_{x_by}=\Gamma^{x_a}_{yx_b}=-\delta^a_b/y$
oraz $\Gamma^y_{yy}=-1/y$; pozostałe symbole są zerowe.
Na przykład pole w kierunku $x$ przenoszone
pionowo musi mieć odpowiednio zmienne
współrzędne: $\nabla_{\partial_y}(y\partial_x)=0$.
Rzeczywiście $\partial_y(y)=1$ oraz
$\Gamma^x_{yx}y=-1$; jego długość w metryce
hiperbolicznej stale wynosi $1$.
\end{przyklad}

\begin{przyklad}[Kula Poincarégo: jawne symbole]
\label{prz:christoffel-dysk}
Metryka~\eqref{eq:hip-kula} ma postać
$g=\lambda(u)^2\sum_i(\dd u^i)^2$,
$\lambda=2/(1-|u|^2)>0$. Indeksy $u_i=u^i$ w tym
przykładzie odnoszą się do współrzędnych euklidesowych;
nie oznaczają obniżenia indeksu metryką hiperboliczną.
Po wstawieniu
$g_{ij}=\lambda^2\delta_{ij}$ i
$g^{ij}=\lambda^{-2}\delta^{ij}$ do
\eqref{eq:christoffel-metryka} otrzymujemy
dla dowolnej metryki takiej postaci
\begin{equation}\label{eq:gamma-konforemne}
 \Gamma^k_{ij}=\delta^k_i\partial_j\log\lambda
 +\delta^k_j\partial_i\log\lambda
 -\delta_{ij}\partial_k\log\lambda.
\end{equation}
Rzeczywiście $\partial_i g_{j\ell}
=2\lambda(\partial_i\lambda)\delta_{j\ell}$;
kontrakcja z $\lambda^{-2}\delta^{k\ell}/2$
daje trzy pokazane wyrazy. Ponieważ
$\partial_k\log\lambda=2u_k/(1-|u|^2)$,
otrzymujemy \emph{jawnie}
\begin{equation}\label{eq:gamma-dysk}
 \boxed{\displaystyle
 \Gamma^k_{ij}=\frac{2}{1-|u|^2}
   \bigl(\delta^k_i u_j+\delta^k_j u_i
                          -\delta_{ij}u_k\bigr).}
\end{equation}
Na przykład w dysku o współrzędnych $(u,v)$,
przy $D=1-u^2-v^2$, niezależne składowe to
\[
 \begin{array}{lll}
 \Gamma^u_{uu}=2u/D,&\Gamma^u_{uv}=2v/D,&
 \Gamma^u_{vv}=-2u/D,\\
 \Gamma^v_{uu}=-2v/D,&\Gamma^v_{uv}=2u/D,&
 \Gamma^v_{vv}=2v/D.
 \end{array}
\]
Wszystkie symbole znikają \emph{w środku} dysku,
ale to nie czyni metryki płaską: pochodne symboli
w~\eqref{eq:riemann-wspolrzedne} z
Rozdziału~\ref{ch:krzywizna} mogą być niezerowe.
Jako kontrolę wzoru weźmy średnicę $u=r$, $v=0$.
Pole $V=\frac{1-r^2}{2}\partial_v$ wzdłuż niej
ma długość $1$ i jest równoległe, gdyż
\[
 \nabla_{\partial_r}V
 =\left(-r+\Gamma^v_{uv}\frac{1-r^2}{2}\right)
             \partial_v
 =\left(-r+\frac{2r}{1-r^2}\frac{1-r^2}{2}\right)
             \partial_v=0.
\]
\end{przyklad}

\begin{przyklad}[Hiperboloida: symbole w mapie przestrzennej]
\label{prz:christoffel-hiperboloida}
Użyjmy współrzędnych $z$ z
\eqref{eq:hiperboloida-wspolrzedne-metryka},
oznaczając $t^2=1+|z|^2$. Także tutaj $z_i=z^i$
oznacza współrzędne w euklidesowej przestrzeni parametrów. Mamy
$g_{ij}=\delta_{ij}-z_i z_j/t^2$
oraz $g^{ij}=\delta^{ij}+z_i z_j$ z
\eqref{eq:hiperboloida-macierz-odwrotna}.
Po zróżniczkowaniu
\[
 \partial_i g_{j\ell}
 =-\frac{\delta_{ij}z_\ell+\delta_{i\ell}z_j}{t^2}
   +\frac{2z_i z_j z_\ell}{t^4},
\]
więc trzy pochodne ze wzoru
\eqref{eq:christoffel-metryka} łączą się w
\[
 \partial_i g_{j\ell}+\partial_j g_{i\ell}
                -\partial_\ell g_{ij}
 =-\frac{2z_\ell}{t^2}
            \left(\delta_{ij}-\frac{z_i z_j}{t^2}\right)
 =-\frac{2z_\ell}{t^2}g_{ij}.
\]
Ponieważ $g^{k\ell}z_\ell=t^2 z^k$, otrzymujemy
\begin{equation}\label{eq:gamma-hiperboloida}
 \boxed{\displaystyle\Gamma^k_{ij}=-z^k g_{ij}
  =-z^k\left(\delta_{ij}-\frac{z_i z_j}{1+|z|^2}\right).}
\end{equation}
Na przykład dla $z(\tau)=(\sinh\tau,0,\ldots,0)$
mamy $g_{11}=1/\cosh^2\tau$ oraz
$\Gamma^1_{11}=-\sinh\tau/\cosh^2\tau$.
Równanie geodezyjnej w pierwszej współrzędnej
przyjmuje zatem postać
$\ddot z^1+\Gamma^1_{11}(\dot z^1)^2
=\sinh\tau-\sinh\tau=0$,
zgodnie ze wzorem~\eqref{eq:exp-hiperboloida}
w rozdziale o geodezyjnych.
\end{przyklad}

Izometrie między modelami z
Przykładu~\ref{prz:trzy-modele-hiperboliczne}
przenoszą te koneksje: wynika to z jednoznaczności
koneksji Leviego-Civity przy zachowaniu zgodności
z metryką i braku torsji. Dokładniej, dla izometrii
$F:(M,g)\to(N,h)$ definiujemy na $M$
\[
 \widehat\nabla_XY=(F^{-1})_*
       \bigl(\nabla^h_{F_*X}(F_*Y)\bigr).
\]
Reguła łańcuchowa daje aksjomaty koneksji. Równości
$F_*[X,Y]=[F_*X,F_*Y]$ i $F^*h=g$ zapewniają zerową
torsję oraz zgodność z $g$. Jednoznaczność daje
$\widehat\nabla=\nabla^g$, czyli
$F_*(\nabla^g_XY)=\nabla^h_{F_*X}(F_*Y)$.
Symboli Christoffela
\emph{nie porównujemy bezpośrednio} między różnymi
mapami: transformują się wraz z pochodnymi zmiany
współrzędnych, podczas gdy sama koneksja jest
tym samym obiektem geometrycznym.

\begin{stwierdzenie}[Dywergencja zapisana przez koneksję]
\label{prop:dywergencja-koneksja}
Dla koneksji Leviego-Civity i pola $X$ na rozmaitości
riemannowskiej tensor $\nabla X$ przyporządkowuje wektorowi
$Y_p$ wektor $(\nabla_YX)_p$. Jego ślad jest dywergencją:
\begin{equation}\label{eq:dywergencja-koneksja}
 \operatorname{div}_gX=\nabla_iX^i
 =\partial_iX^i+\Gamma^i_{ij}X^j
 =\frac1{\sqrt{\det G}}\partial_i
       (\sqrt{\det G}\,X^i).
\end{equation}
Jest to ten sam operator, który wystąpił
w Twierdzeniu Gaussa w Rozdziale~\ref{ch:metryka-riemannowska}.
\end{stwierdzenie}
\begin{proof}
Podstawiamy \eqref{eq:christoffel-metryka}
i korzystamy z symetrii macierzy. Przy sumowaniu
po $i$ składniki
$g^{i\ell}\partial_i g_{j\ell}$
i $-g^{i\ell}\partial_\ell g_{ij}$ są przeciwne
po zamianie nazw indeksów $i,\ell$.
Pozostaje
\[
 \Gamma^i_{ij}=\tfrac12g^{i\ell}\partial_jg_{i\ell}
 =\partial_j\log\sqrt{\det G}.
\]
Ostatnia równość jest wzorem na pochodną
wyznacznika: $\partial_j\det G
=(\det G)\operatorname{tr}(G^{-1}\partial_jG)$.
Reguła iloczynu daje ostatnie wyrażenie
w \eqref{eq:dywergencja-koneksja}, zgodne
z wcześniejszym \eqref{eq:dywergencja-lokalna}.
\end{proof}

\section{Pochodna kowariantna i transport wzdłuż krzywej}
\label{sec:transport-rownolegly}

Pochodna $\nabla_XY$ odpowiada na pytanie:
jak szybko zmienia się $Y$, gdy ruszamy z punktu
w kierunku $X$, \emph{z uwzględnieniem zmiany lokalnej bazy}?
Wzór~\eqref{eq:koneksja-wspolrzedne} dodaje do
pochodnej współrzędnych dokładnie poprawkę
$\Gamma^k_{ij}X^iY^j$.

\begin{definicja}[Pochodna wzdłuż krzywej]
\index{pochodna wzdluz krzywej@Pochodna wzdłuż krzywej}
\label{def:pochodna-wzdluz-krzywej}
Niech $\gamma:I\to M$ będzie gładką krzywą.
Gładkie pole wzdłuż niej to gładkie odwzorowanie
$V:I\to TM$ z $V(t)\in T_{\gamma(t)}M$,
czyli gładki przekrój wiązki $\gamma^*TM$.
Jego pochodną kowariantną oznaczamy
$D V/\dd t$ lub $\nabla_{\dot\gamma}V$.
W mapie, przy $\gamma(t)=(x^i(t))$ oraz
$V(t)=V^k(t)\partial_k|_{\gamma(t)}$, jest to
\begin{equation}\label{eq:pochodna-kowariantna-krzywa}
 \frac{D V}{\dd t}
 =\left(\frac{\dd V^k}{\dd t}
   +\Gamma^k_{ij}(\gamma(t))\dot x^i V^j\right)
    \partial_k\big|_{\gamma(t)}.
\end{equation}
Pole jest \emph{równoległe} wzdłuż $\gamma$,
gdy $DV/\dd t=0$. Rozwiązanie z wartością
$V(a)=v$ oznaczymy przez $V(t)=P_{\gamma,a\to t}v$.
\end{definicja}

Uzasadnijmy niezależność od mapy bez zakładania, że $V$
ma przedłużenie do pola na $M$. Niech $J(t)$ będzie macierzą
$\partial y/\partial x$ w $\gamma(t)$ i niech
$A^k{}_j=\Gamma^k_{ij}(\gamma(t))\dot x^i$.
Współczynniki pola w nowych współrzędnych to
$\widetilde V=JV$. Wzór~\eqref{eq:transformacja-christoffel}
i pochodna $J^{-1}$ dają
\[
 \widetilde A=JAJ^{-1}-\dot J J^{-1},\qquad
 \frac{\dd\widetilde V}{\dd t}+\widetilde A\widetilde V
   =J\left(\frac{\dd V}{\dd t}+AV\right).
\]
Jest to dokładnie prawo transformacji wektora, więc definicja
jest poprawna również przy samoprzecięciach i punktach
z $\dot\gamma=0$. Potrzeba tego argumentu: dla krzywej
stałej $\gamma(t)=p$ pole $V(t)=t v$, $v\ne0$, nie jest
nawet lokalnie ograniczeniem jednego pola na $M$.

Jeżeli $V=Y\circ\gamma$ pochodzi od pola z otoczenia,
reguła łańcuchowa daje $DV/\dd t=\nabla_{\dot\gamma}Y$.
Dla dowolnego pola wzdłuż krzywej mamy też
$D(fV)/\dd t=f'V+f\,DV/\dd t$. Przy zmianie parametru
$t=h(s)$ bezpośrednie podstawienie do wzoru daje
\begin{equation}\label{eq:transport-reparametryzacja}
 \frac{D(V\circ h)}{\dd s}
   =h'(s)\frac{DV}{\dd t}\bigg|_{t=h(s)}.
\end{equation}
W szczególności zachowująca orientację gładka zmiana
parametru nie zmienia transportu między końcami.

\begin{twierdzenie}[Istnienie i własności transportu]
\label{thm:transport-rownolegly}
Dla gładkiej krzywej $\gamma:[a,b]\to M$
i $v\in T_{\gamma(a)}M$ istnieje dokładnie
jedno równoległe pole $V$ wzdłuż $\gamma$
spełniające $V(a)=v$. Transport
$P_{\gamma,a\to b}:T_{\gamma(a)}M\to T_{\gamma(b)}M$
jest liniowym izomorfizmem. Na kolejnych
odcinkach drogi transporty składają się;
po odwróceniu drogi otrzymujemy odwrotność.
Jeżeli koneksja jest zgodna z metryką $g$, to
\begin{equation}\label{eq:transport-izometria}
 g_{\gamma(t)}(P_{a\to t}v,P_{a\to t}w)
   =g_{\gamma(a)}(v,w).
\end{equation}
Dotyczy to również metryk lorentzowskich,
gdzie zachowywana forma nie jest dodatnio określona.
\end{twierdzenie}
\begin{proof}
W jednej mapie warunek $DV/\dd t=0$
jest liniowym układem zwyczajnych równań
różniczkowych
\begin{equation}\label{eq:ode-transport}
 \dot V^k(t)=-A^k{}_j(t)V^j(t),\qquad
 A^k{}_j(t)=\Gamma^k_{ij}(\gamma(t))\dot x^i(t).
\end{equation}
Współczynniki są ciągłe. Twierdzenie o liniowym
układzie ODE daje jednoznaczne rozwiązanie na
każdym zwartym odcinku, na którym krzywa
pozostaje w mapie: rozwiązanie nie może uciec
do nieskończoności w skończonym czasie przy
ciągłej, ograniczonej macierzy $A(t)$. Istotnie, równanie
całkowe i nierówność Grönwalla dają, dla $t\ge t_0$,
\[
 |V(t)|\le |V(t_0)|
       \exp\!\left(\int_{t_0}^t\|A(u)\|\,\dd u\right).
\]
Ograniczoność $A$ i $V$ ogranicza też $\dot V=-AV$,
więc na skończonym końcu przedziału istnieje granica $V$;
lokalne twierdzenie o istnieniu przedłuża rozwiązanie.
Ten sam argument działa wstecz w czasie.
Zwartość $[a,b]$ pozwala podzielić go na
skończenie wiele odcinków, z których każdy
mieści się w jednej mapie. Rozwiązania sklejamy
zgodnie z ich jednoznacznością na przecięciach.

Liniowość układu w $V$ daje liniowość transportu.
Rozwiązując równanie od $b$ ku $a$, dostajemy
odwrotność; jednoznaczność daje też prawo
składania po przecięciu drogi. Wreszcie
zgodność metryczna w lokalnej bazie ma postać
$\partial_i g_{jk}=\Gamma^\ell_{ij}g_{\ell k}
+\Gamma^\ell_{ik}g_{j\ell}$. Różniczkując
$g_{jk}(\gamma(t))V^j(t)W^k(t)$, dostajemy więc,
także bez przedłużania $V,W$ poza krzywą,
\[
 \frac{\dd}{\dd t}g(V,W)
 =g(DV/\dd t,W)+g(V,DW/\dd t).
\]
Gdy $V,W$ są równoległe, prawa strona jest
zerem, więc iloczyn pozostaje stały.
\end{proof}

Dla ciągłej krzywej kawałkami gładkiej, z podziałem na
skończenie wiele gładkich odcinków, definiujemy transport
jako złożenie transportów na tych odcinkach. Jednoznaczność
zapewnia niezależność od zagęszczenia podziału. W ten sposób
prawo składania obejmuje również sklejenia z narożnikiem.

Transport zależy zwykle od \emph{drogi}, a nie
tylko od jej końców. Jeśli koneksja jest zgodna z metryką
i krzywa jest zamknięta, jej transport jest izometrią
$T_pM\to T_pM$;
poniższy przykład pokaże niezerowy obrót
wektora po pełnym okrążeniu.

\begin{przyklad}[Obrót po równoleżniku sfery]
\label{prz:transport-sfera}
Na sferze $S_R^2$, $R>0$, użyjmy lokalnie współrzędnych
$(\theta,\phi)$ i ramki ortonormalnej
\[
 e_\theta=R^{-1}\partial_\theta,
 \qquad e_\phi=(R\sin\theta)^{-1}\partial_\phi.
\]
Idźmy po równoleżniku $\theta=\theta_0$,
$\phi$ od $0$ do $2\pi$, z $0<\theta_0<\pi$.
Cały równoleżnik nie mieści się w jednej mapie kątowej,
ale pola $e_\theta,e_\phi$ sklejają się do gładkiej ramki
na sferze bez biegunów i są $2\pi$-okresowe w $\phi$.
Możemy zatem rozwiązać równanie na całym przedziale.
Wstawiając \eqref{eq:gamma-sfera} do reguły
iloczynu, otrzymujemy
\[
 \nabla_{\partial_\phi}e_\theta
   =\cos\theta_0\,e_\phi,
 \qquad
 \nabla_{\partial_\phi}e_\phi
   =-\cos\theta_0\,e_\theta.
\]
Jeśli $V=A(\phi)e_\theta+B(\phi)e_\phi$,
to równanie $DV/\dd\phi=0$ ma postać
\[
 A'=B\cos\theta_0,
 \qquad B'=-A\cos\theta_0,
 \quad\Longrightarrow\quad
 A+iB=(A(0)+iB(0))e^{-i\phi\cos\theta_0}.
\]
Po powrocie do punktu początkowego ramka
jest znowu ta sama, a wektor obrócił się
o kąt $-2\pi\cos\theta_0$ modulo $2\pi$, mierzony od
$e_\theta$ ku $e_\phi$.
Na równiku kąt jest zerowy. Dla
$\theta_0=\pi/3$ wynosi $-\pi$: wektor
wraca przeciwnie skierowany. Obliczenie
pokazuje różnicę między przenoszeniem
\emph{równoległym na sferze} a utrzymywaniem
strzałki nieruchomej na kartce.
\end{przyklad}

\begin{figure}[htbp]
\centering
\begin{tikzpicture}[>=Latex,scale=1.15]
  % Rzut ortogonalny (X,Y,Z) -> (X, cos(20) Z - sin(20) Y).
  % R=1.75, theta=60 stopni; widoczność: Y cos(20)+Z sin(20)>=0.
  \pgfmathsetmacro{\latradius}{1.75*sin(60)}
  \pgfmathsetmacro{\latcenter}{1.75*cos(60)*cos(20)}
  \pgfmathsetmacro{\latheight}{\latradius*sin(20)}
  \pgfmathsetmacro{\limbangle}{asin(cot(60)*tan(20))}
  \pgfmathsetmacro{\pointheight}{\latcenter-\latheight}
  \shade[inner color=white,outer color=manifoldblue!28] (0,0) circle (1.75);
  \draw[manifoldblue!75!black,thick] (0,0) circle (1.75);
  \draw[manifoldred!65,densely dashed,thick]
    plot[domain=180+\limbangle:360-\limbangle,samples=70,variable=\u]
       ({\latradius*cos(\u)},{\latcenter-\latheight*sin(\u)});
  \draw[manifoldred,very thick]
    plot[domain=-\limbangle:180+\limbangle,samples=70,variable=\u]
       ({\latradius*cos(\u)},{\latcenter-\latheight*sin(\u)});
  \draw[manifoldred,very thick,-{Latex[length=2.5mm]}]
    plot[domain=30:52,samples=15,variable=\u]
       ({\latradius*cos(\u)},{\latcenter-\latheight*sin(\u)});
  \coordinate (p) at (0,\pointheight);
  \draw[manifoldgreen!80!black,ultra thick,->]
    (p) -- ++(0,{-1.5*sin(80)}) node[below] {$v=e_\theta$};
  \draw[manifoldgold!85!black,ultra thick,->]
    (p) -- ++(0,{1.5*sin(80)}) node[above] {$P_\gamma v=-v$};
  \fill[manifoldred] (p) circle (2.2pt) node[below left] {$p$};
  \node[manifoldblue!80!black] at (2.02,-1.15) {$S_R^2$};
  \node[manifoldred!80!black,align=left,anchor=west] at (2.0,.60)
    {$\gamma:\ \theta=\pi/3$\\jeden pełny obieg};
\end{tikzpicture}
\caption{Transport po równoleżniku $\theta=\pi/3$ obraca wektor
o $-\pi$. Dla czytelności początek pętli wybrano w punkcie
$p$ o $\phi=\pi/2$ (zakres parametru przesuwa się wtedy
o $\pi/2$). Strzałki są rzutami przeciwnych wektorów
w tej samej płaszczyźnie $T_pS_R^2$, narysowanymi w tej samej skali.
Przerywana część równoleżnika leży z tyłu sfery.}
\label{fig:transport-sfera}
\end{figure}

\subsection*{Pochodna kowariantna innych pól tensorowych}
\label{subsec:pochodna-kowariantna-tensory}
Pochodna funkcji to $\nabla_X f=X(f)$. Dla
$1$-formy $\alpha$ definiujemy $\nabla_X\alpha$
przez wymóg różniczkowania parowania:
\begin{equation}\label{eq:pochodna-kowariantna-forma}
 (\nabla_X\alpha)(Y)
  =X\bigl(\alpha(Y)\bigr)-\alpha(\nabla_XY).
\end{equation}
Prawa strona jest $C^\infty(M)$-liniowa w $Y$: po
zastąpieniu $Y$ przez $fY$ dwa składniki $X(f)\alpha(Y)$
znoszą się. Określa więc $1$-formę. W argumencie $X$
jest liniowa nad funkcjami, a w $\alpha$ spełnia regułę
Leibniza. We współrzędnych
$(\nabla_i\alpha)_j=\partial_i\alpha_j
-\Gamma^k_{ij}\alpha_k$. Na tensorach
definiujemy $\nabla_X$ przez regułę Leibniza
względem iloczynu tensorowego i zgodność
z kontrakcją (ściąganiem indeksów). Ściślej, dla tensora
$S$ typu $(r,s)$, $1$-form $\alpha_1,\ldots,\alpha_r$
i pól $Y_1,\ldots,Y_s$ wymagamy
\begin{align*}
 (\nabla_XS)(\alpha_1,\ldots,\alpha_r,Y_1,\ldots,Y_s)
 ={}&X\bigl(S(\alpha_1,\ldots,\alpha_r,Y_1,\ldots,Y_s)\bigr)\\
 &-\sum_{a=1}^r S(\alpha_1,\ldots,\nabla_X\alpha_a,\ldots,
                         \alpha_r,Y_1,\ldots,Y_s)\\
 &-\sum_{b=1}^s S(\alpha_1,\ldots,\alpha_r,
                         Y_1,\ldots,\nabla_XY_b,\ldots,Y_s).
\end{align*}
Reguły Leibniza sprawiają, że pochodne mnożników
w każdym argumencie się znoszą. Wynik jest więc tensorem
typu $(r,s)$, określonym bez wyboru bazy. Rozwinięcie
w lokalnych iloczynach tensorowych wektorów bazowych
i form dualnych dowodzi istnienia i jednoznaczności
rozszerzenia: każdy górny indeks wnosi składnik z $+\Gamma$,
a każdy dolny z $-\Gamma$. Na przykład
\[
 (\nabla_i S)_{jk}
 =\partial_iS_{jk}
  -\Gamma^\ell_{ij}S_{\ell k}
  -\Gamma^\ell_{ik}S_{j\ell}.
\]
Dla koneksji zgodnej z metryką dodatkowo $\nabla_i g_{jk}=0$;
jest to właśnie warunek~\eqref{eq:zgodnosc-metryczna}.
Dowolna koneksja afiniczna nie musi go spełniać.
W przypadku koneksji bez torsji pochodna Liego
z Rozdziału~\ref{ch:pola} i pochodna
kowariantna mają związek
\begin{equation}\label{eq:lie-koneksja}
 \Lie_XY=[X,Y]=\nabla_XY-\nabla_YX.
\end{equation}
Pochodna Liego porównuje pola przez przepływ
$X$, a pochodna kowariantna wykorzystuje
wybraną koneksję. Równość jest konsekwencją
zerowej torsji, nie definicją którejś z pochodnych.

\section{Przykłady w metrykach Schwarzschilda i Kerra}
\label{sec:koneksja-schwarzschild-kerr}

W obu czasoprzestrzeniach z Rozdziału~\ref{ch:metryka-riemannowska}
macierz metryki jest nieosobliwa w podanej tam
dziedzinie. Wzór Koszula daje zatem jednoznaczną
koneksję także tutaj. Transport zachowuje
\emph{lorentzowski iloczyn} wektorów;
nie interpretujemy $g(v,v)$ jako kwadratu
euklidesowej długości dla wszystkich wektorów.

\begin{przyklad}[Wybrane symbole Schwarzschilda]
\label{prz:koneksja-schwarzschild}
Niech $f(r)=1-2M/r$, zatem $f'(r)=2M/r^2$.
Dla metryki~\eqref{eq:schwarzschild}
macierz jest diagonalna:
\[
 (g_{\mu\nu})
 =\operatorname{diag}(-f,f^{-1},r^2,r^2\sin^2\theta),
 \qquad g^{rr}=f,\quad g^{tt}=-f^{-1}.
\]
Liczymy kilka przydatnych składników
za pomocą \eqref{eq:christoffel-metryka}:
\begin{align}\label{eq:gamma-schwarzschild}
 \Gamma^t_{tr}=\Gamma^t_{rt}
   &=\frac{f'}{2f}=\frac{M}{r^2f},
 &\Gamma^r_{tt}&=\frac{ff'}2=\frac{Mf}{r^2},\notag\\
 \Gamma^r_{rr}&=-\frac{f'}{2f},
 &\Gamma^r_{\theta\theta}&=-rf,
 &\Gamma^r_{\phi\phi}&=-rf\sin^2\theta,\\
 \Gamma^\theta_{r\theta}=\Gamma^\theta_{\theta r}
   &=\frac1r,
 &\Gamma^\phi_{r\phi}=\Gamma^\phi_{\phi r}&=\frac1r.
 \notag
\end{align}
Na przykład
$\Gamma^r_{tt}=-\tfrac12g^{rr}\partial_rg_{tt}
=-\tfrac12f(-f')=ff'/2$, a
$\Gamma^r_{\theta\theta}
=-\tfrac12 f\partial_r(r^2)=-rf$.
Pozostałe symbole kątowe są identyczne jak
w \eqref{eq:gamma-sfera}; pełny spis nie jest
potrzebny do użycia metody.

Dla drogi radialnej
$\gamma(s)=(t_0,r(s),\theta_0,\phi_0)$
wektor $V(s)=V^\theta(s)\partial_\theta$
jest równoległy wtedy i tylko wtedy, gdy
\[
 \frac{\dd V^\theta}{\dd s}
 +\frac{\dot r}{r}V^\theta=0,
 \qquad\Longleftrightarrow\qquad
 r(s)V^\theta(s)=\text{stała}.
\]
To łatwy test wzoru: $g_{\theta\theta}=r^2$,
więc $g(V,V)=r^2(V^\theta)^2$ jest zachowane.
Chodzi o transport w pełnej czasoprzestrzeni
wzdłuż tej wybranej drogi, nie o wyznaczenie
trajektorii fizycznej.
\end{przyklad}

\begin{przyklad}[Wybrane symbole Kerra]
\label{prz:koneksja-kerr}
Stosujemy $\Sigma=r^2+a^2\cos^2\theta$,
$\Delta=r^2-2Mr+a^2$ oraz współrzędne
Boyera--Lindquista z \eqref{eq:kerr}.
W bloku $(r,\theta)$ metryka i jej odwrotność
są diagonalne:
\[
 g_{rr}=\frac{\Sigma}{\Delta},\quad
 g_{\theta\theta}=\Sigma,\quad
 g^{rr}=\frac{\Delta}{\Sigma},\quad
 g^{\theta\theta}=\frac1\Sigma.
\]
Ponieważ $\partial_r\Sigma=2r$ oraz
$\partial_r\Delta=2(r-M)$, otrzymujemy
\begin{align}\label{eq:gamma-kerr-proste}
 \Gamma^r_{rr}
   &=\tfrac12g^{rr}\partial_rg_{rr}
     =\frac r\Sigma-\frac{r-M}{\Delta},
 & \Gamma^r_{\theta\theta}
   &=-\tfrac12g^{rr}\partial_rg_{\theta\theta}
     =-\frac{r\Delta}{\Sigma},\notag\\
 \Gamma^\theta_{r\theta}
  =\Gamma^\theta_{\theta r}&=\frac r\Sigma,
 &\Gamma^r_{r\theta}
  =\Gamma^r_{\theta r}
  &=-\frac{a^2\sin\theta\cos\theta}{\Sigma}.
\end{align}
Ostatni wzór pochodzi z
$\frac12 g^{rr}\partial_\theta g_{rr}
=\frac{\partial_\theta\Sigma}{2\Sigma}$.
Pokazuje, że współczynnik radialny może
zmieniać się także przy zmianie kąta $\theta$.

Parametr obrotu $a$ wprowadza sprzężenie
między czasem a kierunkiem $\phi$.
Z \eqref{eq:kerr} mamy
$g_{t\phi}=-2Mar\sin^2\theta/\Sigma$.
Metryka nie zależy od $t,\phi$ i nie ma
wyrazów $\dd r\,\dd t$ ani $\dd r\,\dd\phi$,
zatem \eqref{eq:christoffel-metryka} daje
\begin{align}\label{eq:gamma-kerr-mieszany}
 \Gamma^r_{t\phi}=\Gamma^r_{\phi t}
 &=-\frac12\frac\Delta\Sigma\partial_rg_{t\phi}
  =-\frac{Ma\Delta\sin^2\theta
                (r^2-a^2\cos^2\theta)}{\Sigma^3},\notag\\
 \partial_rg_{t\phi}
 &=\frac{2Ma\sin^2\theta
                (r^2-a^2\cos^2\theta)}{\Sigma^2}.
\end{align}
Przy $a=0$ składnik mieszany znika,
$\Sigma=r^2$, $\Delta=r^2f$;
\eqref{eq:gamma-kerr-proste} redukuje się
do radialnych i kątowych wzorów Schwarzschilda.
To sprawdzian rachunku oraz zależności obu metryk.
\end{przyklad}

Nie należy odczytywać pojedynczego symbolu
Christoffela jako niezależnej od współrzędnych
„siły”. Wzór~\eqref{eq:transformacja-christoffel}
zmienia jego wartość po zmianie mapy.
Geometryczny sens mają całe równania
transportu i pochodne kowariantne.

\section{Dlaczego koneksja pojawia się w równaniu ruchu?}
\label{sec:energia-koneksja}

Funkcjonał energii z \eqref{eq:dlugosc-energia}
pozwala rozpoznać krzywe stacjonarne.
W tej sekcji zakładamy $a<b$ i gładkość $\gamma$.
Rozpatrujemy rozmaitość bez brzegu albo krzywą zawartą
w jej wnętrzu, aby dopuszczalne były dowolne małe wariacje.
Dla ustalonych końców i lokalnego zapisu
$\gamma(t)=(x^i(t))$ mamy
\begin{equation}\label{eq:energia-lokalna-koneksja}
 E(\gamma)=\frac12\int_a^b
       g_{ij}(x(t))\dot x^i(t)\dot x^j(t)\,\dd t.
\end{equation}
Niech $x(t,s)$ będzie wariacją z końcami
niezależnymi od $s$ i
$\eta^k(t)=\partial_sx^k(t,s)|_{s=0}$.
Różniczkując pod całką i uwzględniając
symetrię $g_{ij}$, otrzymujemy
\[
 \left.\frac{\dd E}{\dd s}\right|_{s=0}
 =\int_a^b\left(
  \tfrac12\partial_k g_{ij}\,\eta^k\dot x^i\dot x^j
  +g_{kj}\dot\eta^k\dot x^j\right)\dd t.
\]
Całkując drugi składnik przez części,
wyraz brzegowy $[g_{kj}\eta^k\dot x^j]_a^b$
jest zerowy. Dostajemy
\[
 \delta E=\int_a^b\eta^k
 \left(\tfrac12\partial_k g_{ij}\dot x^i\dot x^j
      -\frac{\dd}{\dd t}(g_{kj}\dot x^j)\right)\dd t.
\]
Jeśli $\delta E=0$ dla wszystkich $\eta$
o zwartym nośniku we wnętrzu przedziału,
lemat podstawowy rachunku wariacyjnego
daje dla każdego $k$
\[
 g_{kj}\ddot x^j
 +\partial_i g_{kj}\dot x^i\dot x^j
 -\tfrac12\partial_k g_{ij}\dot x^i\dot x^j=0.
\]
Wyjaśnijmy dwa użyte tu kroki. Każde gładkie $\eta$ o
zwartym nośniku wewnątrz odcinka zawartego w mapie
realizuje wariacja $x(t,s)=x(t,0)+s\eta(t)$ dla małego
$|s|$, sklejona z niezmienioną krzywą poza tym odcinkiem.
Zwartość nośnika zapewnia, że pozostajemy w dziedzinie mapy.
Jeśli zaś ciągły współczynnik przy pewnym $\eta^k$ w całce
byłby dodatni (lub ujemny) w jednym punkcie, miałby ten sam
znak w małym otoczeniu. Nieujemna, niezerowa funkcja gładka
o nośniku w tym otoczeniu dawałaby niezerową całkę.
To dowód potrzebnej wersji lematu podstawowego.

Iloczyn $\dot x^i\dot x^j$ jest symetryczny;
zastępujemy więc $\partial_i g_{kj}$ przez
$\tfrac12(\partial_i g_{kj}+\partial_jg_{ki})$.
Po pomnożeniu przez $g^{\ell k}$ wzór
\eqref{eq:christoffel-metryka} daje
\begin{equation}\label{eq:rownanie-geodezyjnej-zapowiedz}
 \ddot x^\ell+\Gamma^\ell_{ij}\dot x^i\dot x^j=0,
 \qquad\text{czyli}\qquad
 \frac{D\dot\gamma}{\dd t}=0.
\end{equation}
Rachunek jest lokalny: wystarczy stosować takie wariacje
w kolejnych mapach pokrywających krzywą. Odwrotnie,
jeśli $D\dot\gamma/\dd t=0$, to dla dowolnej gładkiej
wariacji o ustalonych końcach całkowanie przez części
na odcinkach zawartych w mapach daje $\delta E=0$.
Wyrazy brzegowe w punktach podziału znoszą się, ponieważ
każdy jest wartością tego samego iloczynu
$g(\partial_s\gamma,\dot\gamma)$ z przeciwnymi znakami.
Stacjonarność jest więc równoważna temu równaniu.

Krzywa stacjonarna energii przenosi więc
swój wektor prędkości równolegle. Takie krzywe
nazywamy \emph{geodezyjnymi} z parametrem
afinicznym. Zgodność z metryką daje
$\frac{\dd}{\dd t}g(\dot\gamma,\dot\gamma)
=2g(D\dot\gamma/\dd t,\dot\gamma)=0$;
w przypadku riemannowskim prędkość ma zatem stałą normę.
Określenie parametru jest istotne: dla $\widetilde\gamma(s)
=\gamma(h(s))$ reguła~\eqref{eq:transport-reparametryzacja} daje
\[
 \frac{D\dot{\widetilde\gamma}}{\dd s}
 =h''(s)\dot\gamma(h(s))
       +(h'(s))^2\frac{D\dot\gamma}{\dd t}\bigg|_{t=h(s)}.
\]
Dla niestałej geodezyjnej $\dot\gamma$ nigdzie nie znika
(jest transportem niezerowego wektora), więc gładka
zmiana parametru zachowuje równanie dokładnie wtedy,
gdy $h(s)=cs+d$ z $c\ne0$.
Dalsze własności geodezyjnych i odwzorowanie
wykładnicze omówimy w następnym rozdziale,
a krzywiznę w dalszej części skryptu. Dla metryki lorentzowskiej
to samo równanie ma sens jako równanie
stacjonarności odpowiedniego działania;
nie twierdzimy, że definiuje odległość
riemannowską na czasoprzestrzeni.

\par\smallskip
{\footnotesize\noindent\emph{Dalsza lektura:}
\href{https://www.math.toronto.edu/mein/teaching/LectureNotes/rieall.pdf}
{Toronto} (koneksja i transport),
\href{https://web.mit.edu/8.962/www/lecnotes/8_962TA-lec-all.pdf}
{MIT} (symbole Schwarzschilda),
\href{https://sites.science.oregonstate.edu/coursewikis/GGR/book/content/kerr}
{Oregon} (metryka Kerra).\par}
